
\documentclass[aspectratio=169]{beamer}

\mode<presentation>
\usetheme{Boadilla}
\definecolor{NTHU1}{RGB}{127,16,132}
\definecolor{NTHU2}{RGB}{228,210,231}
\definecolor{blue}{RGB}{30,90,205}
\definecolor{red}{RGB}{213,94,0}
\definecolor{green}{RGB}{0,128,0}
\definecolor{charcoalgray}{RGB}{108,104,104}
\setbeamercolor{title}{fg=NTHU1}
\setbeamercolor{frametitle}{fg=NTHU1}
\setbeamercolor{block title}{bg=NTHU1, fg=white}
\setbeamercolor{block body}{bg=NTHU2}
\setbeamercolor{structure}{fg=NTHU1}
\setbeamercolor{item projected}{fg=white}
\setbeamercolor{item}{fg=NTHU1}
\setbeamercolor{subitem}{fg=NTHU1}
\setbeamercolor{section in toc}{fg=NTHU1}
\setbeamercolor{description item}{fg=NTHU1}
\setbeamercolor{caption name}{fg=NTHU1}
\setbeamercolor{button}{bg=NTHU1, fg=white}
\setbeamercolor{caption name}{fg=NTHU1}
\usepackage{graphics}
\usepackage{tikz}
\usepackage{amsmath}
\usepackage{bbm}
\usetikzlibrary{decorations.pathreplacing}
\usepackage{geometry}
\usepackage{booktabs}
\usepackage{bookmark}
\usepackage{multirow, makecell}
\usepackage{float}
\usepackage{fancyvrb}
\usepackage{caption}
\captionsetup[figure]{singlelinecheck = false, format= hang, justification=centering}
\captionsetup[subfigure]{singlelinecheck = false, format= hang, justification=raggedright}
\usepackage{subcaption}
\usepackage{adjustbox}
\usepackage{hyperref}
\usepackage{threeparttable}
\usepackage{appendixnumberbeamer} 
\usepackage[T1]{fontenc}
\usepackage[default]{lato} 
\usepackage{smartdiagram}
\smartdiagramset{
  back arrow disabled = true,
  module minimum width=1cm,
  module x sep = 3,
  uniform arrow color=true,
}
\usepackage{xcolor}
\usepackage{colortbl}
	\definecolor{light-gray}{gray}{0.99}

\newenvironment{wideitemize}{\itemize\addtolength{\itemsep}{10pt}}{\enditemize}
\newenvironment{wideenumerate}{\enumerate\addtolength{\itemsep}{10pt}}{\endenumerate}
\newenvironment{widedescription}{\description\addtolength{\itemsep}{10pt}}{\enddescription}
\hypersetup{
colorlinks=true,
linkcolor=NTHU1,
filecolor=green, 
urlcolor=blue,
}
\beamertemplatenavigationsymbolsempty
\setbeamercolor{author in head/foot}{bg=white, fg=NTHU1}
\setbeamercolor{title in head/foot}{bg=white, fg=NTHU1}
\setbeamercolor{date in head/foot}{bg=white, fg=NTHU1}
\setbeamercolor{section in head/foot}{bg=white, fg=NTHU1}
\setbeamercolor{headline}{bg=white}
\setbeamertemplate{footline}{
    \leavevmode%
    \hbox{%
        \begin{beamercolorbox}[wd=.333333\paperwidth,ht=2.25ex,dp=1ex,center]{date in head/foot}%
            \usebeamerfont{date in head/foot}%\insertdate
        \end{beamercolorbox}%
        \begin{beamercolorbox}[wd=.333333\paperwidth,ht=2.25ex,dp=1ex,center]{title in head/foot}%
            \usebeamerfont{title in head/foot}\insertshorttitle
        \end{beamercolorbox}%
        \begin{beamercolorbox}[wd=.333333\paperwidth,ht=2.25ex,dp=1ex,right]{date in head/foot}%
            \usebeamerfont{date in head/foot}\insertframenumber{}\hspace*{2em}
        \end{beamercolorbox}}%
        \vskip0pt%
    }
    %% May need to script the introduction!!!!!!
%\setbeamercolor{page number in head/foot}{fg=NTHU1}
\setbeamertemplate{section in toc}[sections numbered]
\setbeamertemplate{subsection in toc}{\leavevmode\leftskip=3em\rlap{\hskip-1.75em\inserttocsectionnumber.\inserttocsubsectionnumber}
\inserttocsubsection\par}
\setbeamerfont{subsection in toc}{size=\footnotesize}
%\setbeamertemplate{headline}{%
  %\begin{beamercolorbox}[ht=5.5ex]{section in head/foot}
    %\vskip2pt\insertnavigation{0.33\paperwidth}\vskip2pt
  %\end{beamercolorbox}%
%}
\newenvironment{transitionframe}{\setbeamercolor{background canvas}{bg=white}\setbeamertemplate{footline}{} \begin{frame}[noframenumbering]}{\end{frame}}


\makeatletter
\let\@@magyar@captionfix\relax
\makeatother


\title[]{Public Economics and Development} % Change this regularly
\subtitle[]{Week 6: Theoretical and empirical evidence on state capacity and weak states}
\author[Lee]{Seung-hun Lee }
\institute[]{Taipei School of Economics\\ National Tsing Hua University}

\date[]{April 10th, 2026}

\begin{document}
\begin{transitionframe}
\titlepage
\end{transitionframe}


%%% Color slides for section headers: Use for colloquium version (The ones bewteen \iffals and \fi)

\section{Overview setup of the topic}
\begin{frame}
\frametitle{Drastic differences in capabilities of states across different countries} 
Some of the poorest countries in the world are also  `weak states'
\begin{itemize}\setlength\itemsep{3pt}
  \item State capacity: Encompasses ability to support markets, provide public goods, and raise taxes from broad bases 
   \begin{itemize}\setlength\itemsep{3pt}
    \item Support markets: Enforce contracts, protect property rights
    \item Provide public goods: Education, health, but also peace (by monopolizing violence)
    \item Broad base taxes (income tax) are more rule-based, less burdensome for individuals, but require large investments for implementation 
    \item Narrow base taxes (royalties and tariffs) are easy to implement but could be corruptible and extracts particular group excessively 
  \end{itemize}
\item Weak states: States that cannot support basic economic function, raise revenues, deliver essential public services, and provide law and order
\end{itemize} \medskip 
State capacity is widely dispersed, but also clustered
\begin{itemize}\setlength\itemsep{3pt}
\item Weak states are also poorer and mired in internal conflicts
\item Developed countries: High income, functioning institutions, less internal conflicts
\end{itemize}
\end{frame}

\begin{frame}
\frametitle{Huge variation in state capacity across the globe}
\begin{figure}
  \includegraphics[keepaspectratio,width=0.7\textwidth]{fig1.png}
\end{figure}
\end{frame}

\begin{frame}
\frametitle{State capacity clustered with income and functioning of institutions}
\begin{figure}
  \begin{subfigure}{0.49\textwidth}
  \includegraphics[keepaspectratio,width=\textwidth]{fig2_1.png}
\caption{Income and capacities are clustered}
  \end{subfigure}
  \begin{subfigure}{0.49\textwidth}
  \includegraphics[keepaspectratio,width=\textwidth]{fig2_2.png}
  \caption{Capacities and fragilities index}
  \end{subfigure}
\end{figure}
\end{frame}



\begin{frame}
\frametitle{Why care about state capacity, and open questions?}
Why do state capacities matter?
\begin{itemize}\setlength\itemsep{3pt}
\item Capable states can distribute fruits of economic development and prevent violence that stifles the progress
\begin{itemize}
\item Ability to tax, enforce law, and distribute services contributes to growth and peace
\end{itemize}
\item Ultimately relates to economic and institutional development: Many weak states are poor, authoritarian, and in internal conflict
\end{itemize}\medskip

We discuss open questions on state capacity with theories and empirical findings
\begin{itemize}\setlength\itemsep{3pt}
\item What forces explain the formation of capable states?
\item How do weak states perpetuate violence, corruption, and weak growth?
\item What explains the clustering of state institutions, violence, and income?
\end{itemize}
\end{frame}



\begin{transitionframe}
\frametitle{Roadmap for today}
\tableofcontents
\end{transitionframe}



\section{Conceptualizing state capacity}
\subsection{Besley, Persson (2009): The Origins of State Capacity: Property Rights, Taxation and Politics}
\begin{transitionframe}
\begin{center}
  \begin{huge}
    \textcolor{NTHU1}{%
Besley, Persson (2009):\\ The Origins of State Capacity: Property Rights, Taxation and Politics
    }
  \end{huge}
\end{center}
\end{transitionframe}

\begin{frame}
\frametitle{Q: What determines the state's investment in fiscal and legal capacity}
How are institutional capacity support markets and collect taxes developed?
\begin{itemize}\setlength\itemsep{3pt}
\item Historians note the development of capacity in response to external wars (Tilly 1990)
\item Social sciences (economics, especially) typically assume such capacities exist by default
\item In reality, institutional capacities differ starkly
\end{itemize}\medskip
This paper provides theoretical framework of state institutions as endogenous investments
\begin{itemize}\setlength\itemsep{3pt}
\item Individuals come from different groups (ruling vs nonruling, different wealth)
\item State collects taxes, protects property rights, and provide public goods
\item What determines taxation, public goods provision, and support for market activities
\item Why are there differences in investments in these capacities across different settings?
\begin{itemize}\setlength\itemsep{3pt}
\item Focus on high capacities of wealthy, democratic, and parliamentary democracies 
\end{itemize}
\item Provides grounds for studying how institutions evolve during development
\end{itemize}
\end{frame}

\begin{frame}
\frametitle{Setting up the framework: Definitions and layout} 
What makes up state capacity
\begin{itemize}\setlength\itemsep{3pt}
\item Legal capacity: Enforce contracts, regulate markets, and protect property rights
\item Fiscal capacity: Collect taxes to finance public goods and `redistribution'
\end{itemize}\medskip 
Setting the model up
\begin{itemize}
\item There are two time periods $s\in\{1,2\}$ and two groups of individuals  $J\in\{A,B\}$
\item Individuals have linear preferences over public goods and post-tax income
\item State (led by $A$ or $B$) protects property rights, collect taxes in both periods
\item State can also increase its legal and fiscal capacities at $s=1$
\item State maximizes (politically) weighted sum of individual utilities under budget and capacity constraints
\end{itemize}
\end{frame}

\begin{frame}
\frametitle{Setting up the framework: Individual utilities} 
Individual utilities (For given group $J\in\{A,B\}$)
\begin{itemize}\setlength\itemsep{3pt}
\item Gain utility from public good $G_s$ (by factor of $\alpha_s$) and income $Y_s^J$ (taxed by $t_s^J$)
\[
v_s^J=\alpha_s G_s + (1-t_s^J)Y_s^J(p_s^J) 
\]\vspace{-20pt}
\item Income gain is higher with increase in $p_s^J$, enforcement of property rights, $\left(\frac{\partial Y_s^J}{\partial p_s^J} >0\right)$
\item Taxes deduct from income but used to finance public goods
\end{itemize}\medskip 
Differences across groups
\begin{itemize}\setlength\itemsep{3pt}
\item Demographic share: $\beta^J$ of population in group $J\in\{A,B\}$ ($\beta^A+ \beta^B=1$)
\item Different protection, taxation (negative taxes possible), and income 
\item Public good is distributed commonly (no differences in level of public goods)
\item Assume that group $A$ is ruling at $s=1$, with chance of continued ruling in $s=2$ is $\gamma_A$
\end{itemize}\medskip 

\end{frame}

\begin{frame}
\frametitle{Setting up the framework: State objectives and constraints} 
Objective of the state: Politically weighted sum of individual utilities
\begin{itemize}
\item $\bar{\rho}$ and $\underline{\rho}$: Relative political weights on ruling and nonruling groups ($\bar{\rho}\beta_s^A+\underline{\rho}\beta_s^B=1$)
\[
V_s = \alpha_s G_s + \bar{\rho}\beta_s^A(1-t_s^A)Y_s^A(p_s^A)+\underline{\rho}\beta_s^B(1-t_s^B)Y_s^B(p_s^B)
\]\vspace{-20pt}
\begin{itemize}
\item $\bar{\rho}$ captures how $A$ values its own group relative to demographic share (1 or greater)
\item $\underline{\rho}$ captures how $A$ values nonruling groups (less or equal to 1)
\end{itemize}
\item At $s=1$, invest in legal ($\pi_s$) and fiscal capacity ($\tau_s$) at cost of $L(\pi_2-\pi_1)$ and $F(\tau_2-\tau_1)$
\end{itemize} \medskip
Constraints of the state
\begin{itemize}\setlength\itemsep{3pt}
\item Capacity constraints: $p_s^J \leq \pi_s$, $t_s^J \leq \tau_s$
\item Budget Constraints
\begin{itemize}\setlength\itemsep{3pt}
\item $s=1$: Capacity investment included $\to \sum_J \beta^J t_1^J Y_1^J=G_1 + L(\pi_2-\pi_1)+F(\tau_2-\tau_1)$ 
\item $s=2$: No new investment$\to \sum_J \beta^J t_2^J Y_2^J=G_2 $ 
\end{itemize}
\end{itemize}
\end{frame}

\begin{frame}
\frametitle{Policy choices in utilitarian vs politically-motivated settings} 
Policy choices on $p_s^J$ and $t_s^J$ in utiltarian governments ($\bar{\rho}=\underline{\rho}=1$)
\begin{itemize}\setlength\itemsep{3pt}
\item $p_s^J$ increases indirect utilities: Provide to the fullest possible $p_s^J=\pi_s$ 
\item $t_s^J$ works in both ways - incomes are taxed, but used for public goods
\item $\alpha_s\geq1$: Tax to the fullest possible ($t_s^J=\tau_s$), all revenues to $G_s$ and investments 
\item $\alpha_s<1$: $G_s=0$, smallest possible taxes (0 in $s=2$, exact cost of investments in $s=1$)
\end{itemize}\medskip
Politically-motivated objective maximization ($\bar{\rho}>1>\underline{\rho}$)
\begin{itemize}\setlength\itemsep{3pt}
\item No different logic in legal capacity: Higest protection possible
\item Public goods are provided only if ruling class values it sufficiently ($\alpha_s\geq\bar{\rho}$)
\item If $\alpha<\bar{\rho}$:  $G_s=0$, and taxes could be `redistributed' from nonruling to ruling class
\begin{itemize}
  \item Ruling group gains $\bar{\rho}$ from `redistribution', nonruling loses $\underline{\rho}$ from this
  \item With $\bar{\rho}-\underline{\rho}>0$ beneficial to extract (from nonruling group standpoint)
\item Implies that in this setting, people have different interests towards fiscal capacity
\end{itemize}
\end{itemize}
\end{frame}

\begin{frame}
\frametitle{Investing in state capacity} 
Investments in fiscal ($\tau$) and legal ($\pi$) capacities maximizes expected social utility
\begin{itemize}\setlength\itemsep{3pt}
\item Factors: Wealth ($Y_s^J$), Value of public goods ($\alpha_s$), Regime stability ($\gamma_A$), representativeness ($\bar{\rho}-\underline{\rho}$)
\item Legal capacity is desirable in any case
\item When $\alpha_s<\bar{\rho}$, high $\tau$ undesirable for non-ruling group $\to$ Underinvest in capacities
\begin{itemize}
  \item In other cases, no conflict over the desirability of fiscal capacity
\end{itemize}
\end{itemize}\medskip
Key findings (in theory and empirics)
\begin{itemize}\setlength\itemsep{3pt}
\item In most cases, legal and fiscal capacities complement each other
\begin{itemize}
  \item $\pi\Uparrow\to$ More revenues for public goods $\to\tau\Uparrow\to$ More funds to increase protection$\to\pi\Uparrow$
\end{itemize}
\item What factors increase investments in state capacities?
\begin{itemize}
  \item High $Y_s^J$, High $\alpha_s$, High $\gamma_A$, Small $\bar{\rho}-\underline{\rho}$
  \item Mismatch in political power and economic influence $\to$ Fiscal capacity underinvested
\end{itemize}
\end{itemize}
\end{frame}

\begin{frame}
\frametitle{Theoretical prediction matches well with observed correlates} 
\begin{figure}
  \includegraphics[keepaspectratio, width=0.7\textwidth]{bp2009_t1.png}
\end{figure}
\end{frame}

\begin{frame}
\frametitle{Theoretical prediction matches well with observed correlates} 
\begin{figure}
  \includegraphics[keepaspectratio, width=0.7\textwidth]{bp2009_t2.png}
\end{figure}
\end{frame}




\subsection{Dincecco, Katz (2014): State Capacity and Long-Run Economic Performance}
\begin{transitionframe}
\begin{center}
  \begin{huge}
    \textcolor{NTHU1}{%
Dincecco, Katz (2014):\\ State Capacity and Long-Run Economic Performance
    }
  \end{huge}
\end{center}
\end{transitionframe}


\begin{frame}
\frametitle{Q: long-run relationship btw state capacity and economic growth}
Effective states are only a recent phenomenon
\begin{itemize}\setlength\itemsep{3pt}
\item Having enough administrative infrastructure to secure property rights and collect tax
\item Modern states underwent (1) fiscal centralization and (2) limited government
\begin{itemize}\setlength\itemsep{3pt}
\item Uniform tax systems at national level
\item Establishment of parliaments that monitor state expenditures
\end{itemize}
\end{itemize}\medskip
This paper explores relationship between state capacity and economic performance
\begin{itemize}\setlength\itemsep{3pt}
\item Data from European 11 countries from 1650-1913
\item Fiscal centralization contribute to long-run economic growth
\item Higher extraction of revenues explain how both transformations affect growth
\end{itemize}
\end{frame}

\begin{frame}
\frametitle{Political transformation}
Fiscal centralization: From fragmentation to centralized collection
\begin{itemize}\setlength\itemsep{3pt}
\item Central govt competed with local institutionss in tax collection
\item Local institutions were free-riding on other localities
\item National governments standardized tax implementation (vs bargaining place by place) 
\end{itemize}\medskip
Limited government
\begin{itemize}\setlength\itemsep{3pt}
\item Parliaments demanded fiscal oversight to prevent monarchs from wasteful spending
\item Having regular control led to rise in extractive capacity
\end{itemize}\medskip
Economic implications: How transformations contribute to growth
\begin{itemize}\setlength\itemsep{3pt}
\item Provision of productive infrastructure (education, public services etc)
\item Efficient mobilization of resources (by efficient information gathering and transmission)
\end{itemize}
\end{frame}

\begin{frame}
\frametitle{Data and regression}
Data sources
\begin{itemize}\setlength\itemsep{3pt}
\item Apply different timing of fiscal centralization (e.g. ENG: 1066, FRA: 1790, DEN:1903)
\item Define timing for limited government (ENG: 1688, FRA: 1870, DEN: 1848)
\item State capacity: Per capita national govt revenue and military expenditure 
 \item Economic performance: per capita GDP from Maddison (2010) data
 \begin{itemize}
 \item 1600, 1700, then annually from 1820-1913
 \end{itemize}
\end{itemize}\medskip
Identification: Leverage variation in timing of transformations 
\[
\Delta y_{it}=\alpha + \beta_1 C_{it}+\beta_2 L_{it}+\gamma X_{it-1}+\mu_i + \lambda_t+e_{it}
\]\vspace{-20pt}
\begin{itemize}\setlength\itemsep{3pt}
\item $C_{it}$ and $L_{it}$ are fiscal centralization and limited government dummy
\item $X_{it-1}$: External/internal conflicts, population growth, and lags of dependent variables
\item Country and year fixed effects included
\item Potential reverse causality concerns (why lags of dependent variable are in $X_{it-1}$)
\end{itemize}
\end{frame}

\begin{frame}
\frametitle{Growth in economic performance after political transformations} 
\begin{figure}
  \includegraphics[keepaspectratio, width=0.5\textwidth]{dk2014_f1.png}
\end{figure}
\end{frame}

\begin{frame}
\frametitle{Effects are larger for fiscal centralizations} 
\begin{figure}
  \includegraphics[keepaspectratio, width=0.7\textwidth]{dk2014_t1.png}
\end{figure}
\end{frame}

\begin{frame}
\frametitle{Extractive capacity increases explains the findings} 
\begin{figure}
  \includegraphics[keepaspectratio, width=0.65\textwidth]{dk2014_t2.png}
\end{figure}
\end{frame}


\begin{frame}
\frametitle{Takeaways and conclusions}
Key findings
\begin{itemize}\setlength\itemsep{3pt}
\item Fiscal centralization (and limited govt, to some extent) contribute to long-run growth
\item Particularly driven by rise in extractive capacity
\item The findings implies that growth of administrative infrastructure had significant role
\end{itemize}\medskip
Remaining questions
\begin{itemize}\setlength\itemsep{3pt}
\item  Potential role for non-fiscal mechanisms?
\begin{itemize}\setlength\itemsep{3pt}
\item Fiscal centralization removes political power of local elites (e.g. feudal landlords)
\item Political accountability: Prevent wasteful spending and have voice in policy decisions
\end{itemize}
\item How are taxes collected used is also a remaining question
 \begin{itemize}\setlength\itemsep{3pt}
\item Efficient bureaucracy that provides services at efficient costs?
\end{itemize}
\end{itemize}
\end{frame}


\subsection{Acemoglu, De Feo, De Luca (2020): Weak States; Causes and Consequences of Sicilian Mafia}
\begin{transitionframe}
\begin{center}
  \begin{huge}
    \textcolor{NTHU1}{%
Acemoglu, De Feo, De Luca (2020):\\ Weak States; Causes and Consequences of Sicilian Mafia
    }
  \end{huge}
\end{center}
\end{transitionframe}




\begin{frame}
\frametitle{Q: How do illicit organizations emerge and contribute to weak states?}
Criminal organizations are widespread and may play a significant role in economic growth
\begin{itemize}\setlength\itemsep{3pt}
\item Think of Mafias in Italy, drug cartels in Mexico/Colombia, and so on
\item Many of these have filled the void left by weak states
\begin{itemize}\setlength\itemsep{3pt}
\item Some of these groups have provided public services on their own
\item But also collected `taxes' in exchange for `protection'
\end{itemize}
\item Possibly contribute to weakness of the state and economic underdevelopment?
\end{itemize}\medskip
This paper explores the growth and consequence of Sicillian Mafia
\begin{itemize}\setlength\itemsep{3pt}
\item What explains Mafia's growth dating back to 1860s?
\item How does proliferation of Mafia affects economic development and other outcomes?
\begin{itemize}\setlength\itemsep{3pt}
\item Literacy rates and high school completion rates (provision of education)
\item Possible mechanism: Capture (literal and figurative) of politicians and candidates
\end{itemize}
\end{itemize}
\end{frame}



\begin{frame}
\frametitle{Proliferation of Mafia in Sicily}
Sicilian Mafia origins
\begin{itemize}\setlength\itemsep{3pt}
\item Mafia started during the transition from Bourbon Kingdom to Unified Italy in 1860s
\item Unified Italy not fully set up in Sicilian Islands $\to$ Weak state presence
\item Rise of socialist movement (\textit{Fasci}), especially in the severe droughts of 1893
\item This movement led landowners and local politicians turn to Mafia for protection
\end{itemize}\medskip
Impact of the drought of 1893
\begin{itemize}\setlength\itemsep{3pt}
\item Conditions for agriculture were deteriorating even prior to 1893
\item 1893: Drop in crop production $\to$ Fasci movement $\to$ landowners turn to Mafia
\item This pattern is confirmed from data obtained by the authors
\end{itemize}
\end{frame}

\begin{frame}
\frametitle{Mafia and Draught distribution significantly overlap}
\begin{figure}
  \begin{subfigure}{0.48\textwidth}
  \includegraphics[keepaspectratio,width=\textwidth]{add2020_f1_1.png}
\caption{Draughts in Sicily}  
\end{subfigure}
  \begin{subfigure}{0.48\textwidth}
  \includegraphics[keepaspectratio,width=\textwidth]{add2020_f1_2.png}
  \caption{Mafia presence in Sicily}
  \end{subfigure}
\end{figure}
\end{frame}


\begin{frame}
\frametitle{Data and regression}
Database
\begin{itemize}\setlength\itemsep{3pt}
\item Mafia presence: Obtained from Cutrera (1900)
\item Fasci presence and rainfall: Renda (1977) and Eredia (1918), respectively
\item Human capital outcome: Obtained from Italian census (1911,21,31,61,71, 81)
\item Other public goods: Various mortality rates and expenditure from ISTAT
\end{itemize}\medskip
Identification leverages exogenous variation in Mafia presence explained by draughts
\begin{itemize}\setlength\itemsep{3pt}
\item Reduced form: $y_i=\alpha+ \beta\text{Mafia}_i+\gamma X_i+e_i$
\item IV (first-stage): $\text{Mafia}_i=\alpha_1 + \beta_1 \text{rainfall}_i+\gamma_1 X_i+u_i\Rightarrow \widehat{\text{Mafia}}_i=\widehat{\alpha}_1 + \widehat{\beta}_1 \text{rainfall}_i+\widehat{\gamma}_1 X_i$
\item IV (second-stage): $y_i=\alpha+ \beta\widehat{\text{Mafia}}_i+\gamma X_i+e_i$
\item Relevance: $F$ statistics around 10 (not very strong, but not weak either)
\item Exclusion: Droughts only affect outcomes through Mafia presence
\end{itemize}
\end{frame}

\begin{frame}
\frametitle{Significant mid-term effect on human capital} 
\begin{figure}
  \includegraphics[keepaspectratio, width=1\textwidth]{add2020_t1.png}
\end{figure}
\end{frame}

\begin{frame}
\frametitle{Findings robust to other measures of public service provision} 
\begin{figure}
  \includegraphics[keepaspectratio, width=1\textwidth]{add2020_t2.png}
\end{figure}
\end{frame}

\begin{frame}
\frametitle{Long-run negative (but weaker) effects on human capital} 
\begin{figure}
  \includegraphics[keepaspectratio, width=0.7\textwidth]{add2020_t3.png}
\end{figure}
\end{frame}



\section{Reciprocical nature of state capacity}
\subsection{Besley (2020): State Capacity, Reciprocity, and Social Contract}
\begin{transitionframe}
\begin{center}
  \begin{huge}
    \textcolor{NTHU1}{%
Besley (2020):\\ State Capacity, Reciprocity, and Social Contract
    }
  \end{huge}
\end{center}
\end{transitionframe}

\begin{frame}
\frametitle{Q: How does civic culture contribute to fiscal capacity?}
Increased tax intake enabled governments to spend on essential services 
\begin{itemize}\setlength\itemsep{3pt}
\item Not just on national defense, but also health care, education, income transfers etc
\item Switches in public spending requires economic, political, and social changes
\begin{itemize}\setlength\itemsep{3pt}
\item Should develop capacity to tax and raise revenues
\item But also, cultivate civic culture that increases tax \textit{compliance}
\end{itemize}
\end{itemize}\medskip
This paper extends the Besley, Persson (2009) model
\begin{itemize}\setlength\itemsep{3pt}
\item  By incorporating the role of civic culture, reciprocity, and social contract
\item Relative payoffs of civic-minded citizens and materialistic citizens matter
\item Investigates how civic culture can develop in the long-run
\item Ultimately answers how state capacity can persist with interactions with citizens
\end{itemize}
\end{frame}

\begin{frame}
\frametitle{Motivating findings: Different attitudes to tax compliance} 
\begin{figure}
\includegraphics[keepaspectratio, width=0.7\textwidth]{b2020_t1.png}
\end{figure}
\begin{footnotesize}
  Source: Individual responses from World Value  Survey, regressing compliance on individual characteristics, country and wave dummies. Income is displayed from lowest to highest
\end{footnotesize}
\end{frame}

\begin{frame}
\frametitle{Model setup} 
Basic structure of the model
\begin{itemize}\setlength\itemsep{3pt}
\item Share $e$ of elites who decides public goods and transfers, and non-elites
\item Among non-elites, $\mu_s$ of civic-minded citizens and $1-\mu_s$ of materialists
\begin{itemize}\setlength\itemsep{3pt}
\item All subject to taxes, with possible noncompliance ($n$), and receive utility from public goods 
\item Civic-minded: Gains if tax used on public goods ($G$), loss if used on transfers to elites ($B$)
\item Noncompliance: Incurs cost $C(n)$, which rises with fiscal capacity $\tau$ (easily detected)
\item Individual Utility: $\lambda>0$ for civic-minded citizens, $\lambda=0$ for materialists\vspace{-5pt}
\[\alpha G + b + w(1-(1-n)[t-\lambda\{G-eB\}] -\tau C(n))\] 
\end{itemize}\vspace{-10pt}
\item Revenues must be spent on public goods $G$, transfers to elites ($B$) and non-elites ($b$)
\begin{itemize}\setlength\itemsep{3pt}
\item For every $B$, share $0\leq\sigma\leq1$ must be spent on $b$. Thus $b=\sigma B$  (Higher $\sigma$ $\to$ cohesiveness)
\item Budget: $eB+(1-e)\sigma B = T-G $
\end{itemize}
\item Fiscal capacity is determined by physical investments but also `tax morale'
\begin{itemize}\setlength\itemsep{3pt}
  \item For materialists, only physical investments matter
\item For civic-minded person, public good$\uparrow \to$ tax morale $\uparrow$ (transfer$\uparrow\to$ tax morale $\downarrow$)  
\end{itemize}
\end{itemize}
\end{frame}

\begin{frame}
\frametitle{Optimal policy choice}  
Optimal policy choice
\begin{itemize}\setlength\itemsep{3pt}
\item Elites choose $t$, $G$, and $B$ to maximize the sum of individual utilities
\item Subject to cohesiveness of institutions ($\sigma$), fiscal capacity ($\tau$) and civic culture ($\mu$)
\item Whether revenues are spent on $G$ or $B$ depends on cohesiveness and civic culture
\begin{itemize}
\item Public goods are more likely to be provided in scenario where civic culture ($\mu\Uparrow$) and cohesiveness ($\sigma\Uparrow$) are dominant (and also raise fiscal capacity)
\item This comes from the fact that civic-minded persons are more likely to comply if revenues are used for public goods (and vice versa)
\item Furthermore, high cohesiveness makes transfers to elites expensive $\to$ make public goods relatively cheaper
\end{itemize}
\item Over time, additional public goods due to $(\mu, \sigma) \Uparrow$ $\Rightarrow$ benefits of being civic-minded $\Uparrow$
\begin{itemize}
\item Conversely, narrow range of public goods provision from $(\mu, \sigma) \Downarrow$ $\Rightarrow$ civic-mindedness $\Downarrow$
\end{itemize}
\item Initial state of civicness, cohesiveness + strength of common interest ($\alpha$) matters
\end{itemize}\medskip
\end{frame}

\begin{frame}
\frametitle{Core implications of the model} 
Cohesiveness of institutions and civic culture is complementary
\begin{itemize}\setlength\itemsep{3pt}
\item Cohesiveness of institutions are strengthened when more civic-minded persons induce high public goods over transfers to elites
\item People are more likely to be civic minded if institutional setups make transfers to elites costly (through cohesiveness)
\end{itemize}\medskip
These determine how common interests develop
\begin{itemize}\setlength\itemsep{3pt}
\item In a less cohesive, weakly civic-minded setting, external threat may not significantly increase common interest
\end{itemize}\medskip
Civic-mindedness and cohesivess can be fostered in more representative political setup
\begin{itemize}\setlength\itemsep{3pt}
\item Elites that more closely represents non-elites more likely to sustain power
\item Further increases returns for being civic-minded
\end{itemize}\medskip
\end{frame}

\begin{frame}
\frametitle{Potential extensions} 
High dependence on outside sources of funds
\begin{itemize}\setlength\itemsep{3pt}
\item Outside sources of funding reduces need for own-source tax revenues
\item This reduces importance of tax compliance and the role of civic culture
\item Sudden, exogenous changes to outside funds $\to$ Challenges in raising fiscal capacity 
\end{itemize}\medskip
Can elites determine civic-mindedness ($\mu$)?
\begin{itemize}\setlength\itemsep{3pt}
\item We assumed that $\mu$ is exogenously determined
\item But these can be potentially determined by elites
\item Room for extension: Making tax avoidance costly, Benefits of public services over time
\end{itemize}\medskip
\end{frame}

\begin{frame}
\frametitle{Takeaways: How these topics can be empirically studied} 
How can fiscal capacities be built 
\begin{itemize}\setlength\itemsep{3pt}
\item Evolution towards broad-based system (Gordon, Li 2009; Jensen 2022)
\item Efforts to gather more accurate information in developing countries (Balan et al 2021)
\end{itemize}\medskip
What increases tax morale and compliance?
\begin{itemize}\setlength\itemsep{3pt}
\item Increase benefits of compliance (tax breaks for formal firms) - (Benhassine 2018)
\item Tax perceptions and electoral outcomes (Christensen, Garfias 2021)
\end{itemize}\medskip
Other factors that contribute to state capacities
\begin{itemize}\setlength\itemsep{3pt}
\item Role of bureaucrats: Well-motivated, well-trained bureaucrats matter for effective public goods provision
\item Active studies on how to recruit and monitor them (Dal Bó et al 2013)
\end{itemize}
\end{frame}



\subsection{Weigel (2020): The Participation Dividend of Taxation: How Citizens in Congo Engage More With The State When it Tries to Tax Them }
\begin{transitionframe}
\begin{center}
  \begin{huge}
    \textcolor{NTHU1}{%
Weigel (2020):\\ The Participation Dividend of Taxation: How Citizens in Congo Engage More With The State When it Tries to Tax Them 
    }
  \end{huge}
\end{center}
\end{transitionframe}

\begin{frame}
\frametitle{Q: Do taxation increase civic engagement in fragile states?}
Participation dividend of taxation
\begin{itemize}\setlength\itemsep{3pt}
\item Taxes enable state to provide public goods and services
\item It is said to also stimulate political engagement: Taxation as a social contract
\begin{itemize}
\item Traditionally: Citizens delay paying taxes until rulers made concessions
\end{itemize}
\item Does this form of bargaining also appear in fragile states?
\end{itemize}\medskip
This paper uses randomized rollout of tax campaigns to see if taxes raises participation
\begin{itemize}\setlength\itemsep{3pt}
\item Context: Democratic Republic of Congo
\item Rollout of tax campaign meant to increase citizens registering to pay taxes
\item Active visit of tax collectors in treated towns vs. Declaration system in control group
\item Participation \& perceptions that government are responsible for public goods $\Uparrow$
\end{itemize}
\end{frame}


\begin{frame}
\frametitle{Institutional context of DRC}
Kangara, DRC
\begin{itemize}\setlength\itemsep{3pt}
\item DRC is one of the poorest country in the world (188/200 in terms of tax-GDP ratio)
\item Tax revenue per person at Kangara: \$0.23 
\item Very few people pay taxes to begin with: 6\% of population have ever paid taxes
\end{itemize}\medskip
Tax collection and political participation
\begin{itemize}\setlength\itemsep{3pt}
\item Systematic tax collection began in 2016, induced by redistricting of the region
\item Citizens have limited avenue for political participation
\item Neighborhood meetings, letter of complaint are main channels
\end{itemize}
\end{frame}

\begin{frame}
\frametitle{Research design}
Randomized control treatment
\begin{itemize}\setlength\itemsep{3pt}
\item Door-to-door property tax collection compaign in Apr-Dec 2016
\begin{itemize}
  \item Property registration and tax solicitation
\end{itemize}
\item Treated neighborhoods (253) received this in 2016, (control (178): mid-2017)
\item Stratified randomization based on location and population
\end{itemize}\medskip
Estimation leverages random assignment to treatment controlling for strata
\[
y_{ijk}=\beta I_{jk} + \alpha_k + X_{ijk}\gamma + X_{jk}\phi+e_{ijk}
\]\vspace{-20pt}
\begin{itemize}\setlength\itemsep{3pt}
\item $i$: individual, $j$: neighborhood, $k$: strata
\item Control for individual and neighborhood characteristics ($X_{ijk}$ and $X_{jk}$)
\item Control \& treated groups balanced on observables and tax/political outcome variables
\end{itemize}
\end{frame}


\begin{frame}
\frametitle{Tax compliance rises following treatment} 
\begin{figure}
  \includegraphics[keepaspectratio, width=0.9\textwidth]{w2020_t1.png}
\end{figure}
\end{frame}

\begin{frame}
\frametitle{Citizens' political participation increases} 
\begin{figure}
  \includegraphics[keepaspectratio, width=1\textwidth]{w2020_t2.png}
\end{figure}
\end{frame}

\begin{frame}
\frametitle{Citizen's perception of government responsibility rises} 
\begin{figure}
  \includegraphics[keepaspectratio, width=0.75\textwidth]{w2020_t3.png}
\end{figure}
\end{frame}


\begin{frame}
\frametitle{Takeaways and remaining questions}
Effects of tax collection campaign
\begin{itemize}\setlength\itemsep{3pt}
\item Higher tax compliance (registered properties, willingess to pay taxes)
\item Higher demand for political participation (esp. town halls)
\item Higher recognition of state responsibility and expectation on state's roles
\end{itemize}\medskip
Remaining questions
\begin{itemize}\setlength\itemsep{3pt}
\item Political participation follows tax collection
\begin{itemize}\setlength\itemsep{3pt}
  \item Governments willing and able to respond to demands will proceeds as is
  \item Otherwise, may attempt other ways to raise revenue (e.g. seignorage, tariffs)
\end{itemize}
\item Other ways to increase political particpation and civic engagement?
\begin{itemize}\setlength\itemsep{3pt}
  \item What about people who do not pay taxes (i.e. most people in developing countries)
  \item How do demand from citizens shape politician behavior?
\end{itemize}
\end{itemize}
\end{frame}

\begin{frame}
\frametitle{Unanswered questions in differences in taxation across countries}
Effectiveness of the state is not a given
\begin{itemize}\setlength\itemsep{3pt}
\item Many states are unable to mobilize sufficient resources
\item Which affects their ability to provide essential public services
\item Also have welfare consequences for the citizens
\item What ways to increase capacity?
\begin{itemize}\setlength\itemsep{3pt}
\item Developing countries differ in underlying institutional conditions
\item Many have underdeveloped financial systems and income registries and bureaucracies
\end{itemize}
\end{itemize}\medskip
Raising extractive capacities must be balanced with citizen demand
\begin{itemize}\setlength\itemsep{3pt}
\item Collect taxes $\to$ citizens demand essential services and state responsibility
\item States unable to meet these demands induce tax non-compliance 
\item So how do institutional trust and state capacities interact?
\end{itemize}\end{frame}
 
%%%%%%%%%%%
\end{document}
